\clearpage
\fancyhead[L]{\small IRIS.}
\SourceLeaf{479}
\begin{center}
\rule{0.55\linewidth}{0.4pt}\\[2pt]
Astronomische Nachrichten. Band XXVI. Nr. 611. S. 173. 174. 16. Sept. 1847.\\[2pt]
\rule{0.55\linewidth}{0.4pt}
\end{center}
\begin{center}[Beobachtungen am Meridiankreis von \textsc{Gauss}.]\end{center}

\begin{center}
\small
\setlength{\tabcolsep}{3pt}
\begin{tabular}{lrll}
Iris: 1847 Aug. 21. & $9^{\rm u}\ 52^{\rm m}\ 54^{\rm s}1$ Gött. m. Z. & $297^\circ\ 47'\ 20''8$ & $-13^\circ\ 42'\ 20''4$ \\
\end{tabular}
\end{center}

Die Durchgangszeiten an den einzelnen Fäden stimmten nicht besonders gut überein, weil das Gestirn wegen vorüberziehender Wolken schwer zu sehen war. Unmittelbar darauf breiteten sich die Wolken noch weiter aus, sodass die Kulmination von $\alpha$ Steinbock nicht beobachtet werden konnte.

\begin{center}
\small
\setlength{\tabcolsep}{4pt}
\begin{tabular}{lrrr}
Neptun 1847 Aug. 17. & $12^{\rm u}\ 22^{\rm m}\ 55^{\rm s}2$ & $331^\circ\ 27'\ 14''1$ & $-12^\circ\ 21'\ 5''0$ \\
\end{tabular}
\end{center}

